\section{Source-to-Event Expansion}\label{sec:compression}

A major objective of Motivo is to provide a highly expressive interface for music composition,
reducing the manual effort required to construct complex arrangements.
We evaluate this expressiveness by analyzing the ``compression ratio''---the difference between the amount of
code written
by the composer and the number of low-level MIDI events generated by the compiler.

\subsection{Algorithmic Conciseness}\label{detail-subsec:algorithmic-conciseness}

Traditional MIDI sequencers and digital audio workstations require the composer to manually place every note on a grid.
In contrast, Motivo allows composers to define structural rules and relationships.

Consider the following drum pattern from the \texttt{example.motivo} file:
\begin{lstlisting}[language=Motivo, caption={A generative drum pattern using control flow.},label={lst:evaluation-drum-pattern}]
track percussion using drums {
    pattern groove() {
        for (let i = 0; i < 16; i = i + 1) {
            play hihat from i;
            if (i % 8 == 0) { play kick from i; }
            if (i % 8 == 4) { play snare from i; }
        }
    }
    loop (3) { play groove(); }
}
\end{lstlisting}

This 10-line block of code encapsulates a complex rhythm.
The \texttt{for} loop defines the underlying hi-hat grid, while the modulo-based \texttt{if} conditions procedurally
generate the kick and snare syncopation.
The \texttt{loop (3)} statement then instantiates this 16-beat pattern three times across the timeline.

\subsection{The Compression Ratio}\label{detail-subsec:the-compression-ratio}

To achieve the same result in a traditional MIDI editor, the composer would have to manually draw 48 hi-hat notes, 6
kick drum notes, and 6 snare drum notes---a total of 60 individual note events, which translates to 120 low-level MIDI
messages (\texttt{Note-On} and \texttt{Note-Off}).

Motivo abstracts these 120 MIDI messages into just a handful of declarative rules.
The resulting ``compression ratio'' is exceedingly high.
As the composition grows---for example, looping the groove 32 times instead of 3---the size of Motivo source
code remains
completely unchanged, while the output scales to hundreds or thousands of events.

This expressiveness validates the core thesis: by shifting from a linear, event-based representation to a structural,
programming-based representation, the manual overhead of repetitive event placement can be reduced.
The composer can focus on higher-level musical structure while the lowering engine handles repetitive
temporal event placement.

The main measured expressiveness result is source-to-event expansion.
The 20,000-event single-loop program uses 7 lines of source.
The chord-loop program also uses 7 lines and generates 20,000 note events because each loop iteration emits a
four-note chord.
The 18-line drum groove generates 10,000 note events by combining a 16-step bar pattern with a 500-iteration loop.
These examples demonstrate the intended effect of the language: repeated symbolic structure remains compact in source
form while the compiler expands it into explicit event data.

\begin{figure}[htbp]
  \centering
  \begin{tikzpicture}
    \begin{axis}[
        ybar,
        width=0.95\textwidth,
        height=7cm,
        ymin=0,
        ylabel={Note events generated per source line},
        symbolic x coords={100,1000,10000,20000,chord,drums,example},
        xtick=data,
        x tick label style={rotate=35, anchor=east, font=\scriptsize},
        nodes near coords,
        nodes near coords style={font=\scriptsize, rotate=90, anchor=west},
        bar width=12pt,
        enlarge x limits=0.08,
        ymajorgrids=true,
        grid style={draw=motivoInk!12}
      ]
      \addplot[fill=motivoGreen!70, draw=motivoGreen!90] coordinates {
        (100,14.3)
        (1000,142.9)
        (10000,1428.6)
        (20000,2857.1)
        (chord,2857.1)
        (drums,555.6)
        (example,2.4)
      };
    \end{axis}
  \end{tikzpicture}
  \caption{Source-to-event expansion in the benchmark corpus.}\label{fig:compression-chart}
\end{figure}

This metric should not be overinterpreted.
More generated events per source line is useful when the source expresses meaningful repetition, but it can also produce
uninteresting repetition.
For a bachelor thesis, the metric is still valuable because it connects the language design to a measurable output.
The compiler removes repetitive manual event placement from the source program and makes the expansion explicit in the
MIDI artifact.
